
Large language model evaluation frameworks assess trustworthiness across multiple dimensions---truthfulness, robustness, fairness, safety, privacy---yet report results independently. TrustLLM provides per-dimension accuracy scores; MMTrustEval generates model cards highlighting dimension-specific strengths and weaknesses. This independent evaluation design reflects an implicit assumption: trustworthiness dimensions are uncorrelated. A model achieving 92\% truthfulness accuracy and 89\% fairness accuracy on separate test sets is presumed to exhibit independent performance across these criteria.

Deployment failures challenge this assumption. In safety-critical applications, instances failing multiple trustworthiness criteria simultaneously---compound failures---pose greater risk than single-dimension failures. A medical diagnosis LLM might fail both truthfulness and fairness for the same minority-group patients, yet independent evaluation provides no visibility into this co-occurrence pattern. If truthfulness failures correlate with fairness failures, the actual compound failure rate exceeds what independent scores predict. Current benchmarks do not measure this correlation structure.

The gap matters for model selection. Deployment teams select models based on aggregate scores that mask correlated vulnerabilities. Without coupling awareness, a model with 90\% truthfulness and 90\% fairness might exhibit 81\% joint performance under independence (0.9 $\times$ 0.9) or substantially lower performance under strong negative coupling. Understanding coupling patterns is necessary for estimating compound failure rates in production.

Two competing hypotheses exist in the literature. The broad independence hypothesis, implicit in TrustLLM and MMTrustEval framework design, predicts zero coupling---dimensions reflect orthogonal model capabilities evaluated by distinct test suites. The broad coupling hypothesis, suggested by Li and Li (2024) documenting robustness-fairness trade-offs, predicts pervasive co-occurrence from shared architectural bottlenecks where stress causes cascading failures across many dimension pairs. Neither hypothesis has been empirically tested across all pairwise combinations of trustworthiness dimensions.

Existing multi-dimensional evaluation frameworks lack instance-level binary labels across dimension pairs required to measure coupling. TrustLLM reports per-dimension scores but not co-occurrence statistics showing which instances fail on multiple dimensions. Behavioral detection methods \citep{Zheng2024} validate that trustworthiness properties are observable from model outputs without internal state access, supporting feasibility of coupling measurement from API-accessible evaluation data.

This study demonstrates methodology for measuring coupling breadth and model-specificity in LLM trustworthiness dimensions. The core finding from synthetic validation: coupling is detectable when present (phi 0.33-0.40, p < 1e-13) but sparse---limited to 2 dominant dimension pairs, not pervasive across all 10 possible pairs from 5 dimensions. Truthfulness-robustness coupling (phi 0.36-0.40) appears consistently across three simulated model profiles (GPT-4, Claude-3, Llama-3), while fairness-safety coupling (phi 0.33-0.40) shows comparable strength in synthetic data. No model exhibits coupling in three or more pairs after multiple comparison correction (Bonferroni alpha = 0.01/30 $\approx$ 0.00033), establishing sparsity as a detectable property.

This sparsity persists under difficulty control. Partial correlation analysis yields partial phi values of 0.36-0.56 with effect size retention of 86-161\%. In five of six model-pair combinations in synthetic data, controlling for difficulty strengthens rather than weakens coupling---a suppressor effect where difficulty appears orthogonal to dimension-specific vulnerabilities by design. Quartile stratification shows coupling persistence across difficulty levels (phi $\geq$ 0.25 in 3-4 quartiles per pair), complementing partial correlation findings.

Models show observationally distinct coupling profiles in synthetic data---GPT-4 exhibits a truthfulness-robustness chain, Claude-3 exhibits a fairness-safety-privacy cluster, Llama-3 exhibits minimal coupling---though statistical confirmation via Mantel test remains inconclusive (r < 0.7 criterion met, p > 0.0167 threshold not met). Sample size (n=100 per dimension) is insufficient for Mantel test statistical power; n $\geq$ 500 is required per Legendre and Fortin (2010) power analysis.

All experiments used synthetic coupling data designed to emulate benchmark structure. This validates measurement methodology capabilities---phi coefficient analysis, partial correlation for difficulty control, Mantel test for model comparison---but defers real-world coupling characterization to production benchmark evaluation. Coupling patterns shown demonstrate detectability if coupling exists in real LLMs, not that these specific magnitudes exist in practice. MultiTrust and TrustLLM benchmarks remain access-gated; real model API evaluations were not conducted in this proof-of-concept phase.

Contributions:

\begin{enumerate}
\item Coupling breadth measurement framework via phi coefficient analysis across 5 dimensions and 3 models, establishing coupling as detectable from synthetic behavioral outputs when designed into data (6 significant pairs at phi 0.33-0.40, p < 1e-13).
\end{enumerate}

\begin{enumerate}
\item Difficulty-independence validation framework via dual methods (partial correlation and quartile stratification), demonstrating methodology capability to isolate designed coupling from confounding by instance difficulty (partial phi 0.36-0.56, retention 86-161\%).
\end{enumerate}

\begin{enumerate}
\item Sparsity detection capability via multi-pair analysis with Bonferroni correction, demonstrating methodology can distinguish sparse from broad coupling patterns (0 models meet $\geq 3$ pairs threshold in synthetic data).
\end{enumerate}

\begin{enumerate}
\item Model-specific profile observation via Mantel test, providing observational evidence that methodology could detect distinct coupling patterns given adequate samples (qualitative patterns observable, statistical confirmation underpowered at n=100).
\end{enumerate}

If validated on real benchmarks, these capabilities would enable reframing multi-dimensional trustworthiness evaluation. Rather than assuming universal independence or universal coupling, the methodology could reveal sparse coupling---architecturally independent dimensions by default, with coupling emerging where specific mechanisms overlap. Benchmarks would report coupling statistics alongside per-dimension scores; model selection for deployments requiring multiple trustworthiness guarantees could exploit known coupling patterns.
